% Authoritative LaTeX manuscript; Markdown is a generated compatibility projection.

% <!-- part-title: Orientation and Reference -->

% <!-- chapter-title: Author’s Note -->

% <!-- frontmatter-order: -2 -->

\chapter{Author’s Note}

As a professor, I direct a research lab by identifying promising directions and delegating their detailed investigation to my assistants. I can give those students a broad sense of my goals through meetings and whiteboard drawings, and they carry them out.

% Evidence: E01 | source: book/data/local-evidence/evidentiary-basis-11202026/integration/records/E01.md | kind: author-account
When I began working with GenAI, I treated it much the same way. The agent confidently translated my strategic orders into local implementation and tactical choices. But after about 200,000 lines, I found that although my ability to issue orders was unchanged (if anything, my ambition had grown!), the agents’ ability to carry them out appeared to be decaying. The good vibes had turned bad: how could an agent implement a complicated request one day, yet struggle to correct a simple defect the next? As I examined the software more carefully, I found that it looked and sounded like my vision, but was not. It was a facsimile of my goals, akin to a child’s finger-painted version of the Mona Lisa or a first-year engineering student’s popsicle-stick bridge. My sense of progress had been a mirage. Claude had not fully understood all my intentions, and had not perfectly implemented those it had.

% Evidence: E01 | source: book/data/local-evidence/evidentiary-basis-11202026/integration/records/E01.md | kind: author-account
This book describes what I did next. Firm in the belief that there must be some way to get agents to build reliable systems, I turned to my Professor’s Bookshelf. I revisited the engineering ideas that had shaped my own practice: abstraction, modularity, specifications, and the separation of design from implementation. I began experimenting with different ways of conveying my intentions to agents and different ways of checking that the resulting software honored them. Eventually, it worked.

I share what I learned in the tradition of books on software engineering methodologies. My own understanding of how to engineer was shaped by Dijkstra, Ousterhout, McConnell, Meyer, Uncle Bob, Martin Fowler, Kent Beck, the Gang of Four, and others. Their books taught me how to reason about software, organize its construction, and recognize good engineering. But none of them trained me to delegate implementation to an eager and capable subordinate who could produce extraordinary amounts of code without necessarily understanding what I wanted.

This book teaches the engineering methodology I developed in response. I explain the principles, show how to apply them, and use problems I encountered in DocAble to illustrate the engineering decisions involved: what needed to hold, what I built to achieve it, and what the evidence showed.

It is an intimate thing, to share one's engineering with another. I do not pretend that every decision I made was correct, or that I would make the same decisions again. Some were compromises, some were mistakes, and others made sense only under the circumstances in which they arose. I have tried to show the machinery as it was built, including its limitations, in the hope that you will find something useful for your own engineering.

% <!-- letter-signature -->

James C. Davis\\
West Lafayette, IN\\
2026
